\section{Implementation and Evaluation}
\label{sec:implementation-direction}

Elephant Agent is implemented around the clean Understanding System boundary. The
important implementation property is not that every old object is preserved; it
is that each durable behavior has a single owner and a single correction path.

\subsection{Runtime Records}

The runtime separates Episodes, Loops, and Steps from durable understanding.
Episodes mark open runtime windows. Loops record interaction rounds. Steps
record atomic events for audit and replay. The Personal Model stores active
facts and questions. Elephant State stores the elephant identity plus one natural-language
current context note.

This separation lets the system compact without losing causal structure. A long
tool sequence can be summarized at the Episode level while preserving critical
Step links. A Personal Model fact can be corrected without mutating the
transcript or treating support material as truth.

\subsection{Foreground Update Path}

All foreground durable understanding changes go through
\texttt{tool.personal\_model.update}. The tool writes one lens/topic fact at a
time and supports four actions: remember, correct, forget, and dispute. Init
answers, dashboard corrections, question answers, and chat-time user
corrections all converge on the same path.

\subsection{Claim-Aware Recall and Prompt Projection}

Current-turn retrieval is contextual recall, one layer of the broader memory
architecture. It is injected as \texttt{Current-turn recall support:} when useful
and omitted otherwise. Stable prompt projection contains active
facts only, grouped by Identity, World, Pulse, and Journey, plus the Elephant
identity and Episode opening resume snapshot.

The implementation treats Personal Model search as claim-aware retrieval. The
search tool supports \texttt{topic}, \texttt{query\_variants}, and
\texttt{mode=auto|inventory}. It merges fielded topic/fact/support
matches, Unicode lexical and CJK n-gram signals, weak fuzzy matches, semantic
candidates, and lightweight polarity agreement or conflict. Results are
thresholded into \texttt{strong\_match}, \texttt{weak\_match}, or
\texttt{no\_match}; diagnostics can expose per-fact signals for evaluation but
are not prompt truth.

Topics are represented as facets: stable keys under the four lenses that attach
each claim to one slice of understanding without adding a second profile schema.

\texttt{tool.conversation.search} provides the complementary evidence path. A
small deterministic query planner strips recall operators from the topic core
and detects recency, current-state, historical, recap, and verification intent
across English and CJK phrasing. \texttt{mode=discover} requires a time range
and returns candidate windows; \texttt{mode=recall} retrieves Step or Episode
evidence from a selected window. The searcher first attempts the durable
semantic index with vector, BM25, exact keyword, token-coverage, and n-gram
signals, then falls back to Step scanning when the index is cold. Final ranking
keeps semantic relevance primary and uses temporal intent only as a small
recency or historical tie-breaker.

\subsection{Operational Evaluation}

Elephant Agent is evaluated through longitudinal scenarios rather than single-turn prompt
quality alone:

\begingroup
\small
\setlength{\LTpre}{0.55em}
\setlength{\LTpost}{0.55em}
\begin{longtable}{@{}p{0.22\linewidth}p{0.72\linewidth}@{}}
\toprule
\textbf{Scenario} & \textbf{Expected behavior} \\
\midrule
\endfirsthead
\toprule
\textbf{Scenario} & \textbf{Expected behavior} \\
\midrule
\endhead
Wake recovery & Resume an episode using the Elephant context note, active facts,
and relevant recall. \\
Claim correction & Replace an active fact through \texttt{correct}, retire the
old fact, and project only the new fact next turn. \\
Forget & Retire a fact through \texttt{forget} so it no longer shapes replies. \\
Question answer & Mark the question answered and create the resulting fact
with source Episode provenance. \\
Why inspection & Show source Steps behind a fact without promoting those Steps
to prompt truth. \\
Search no-match & Return no active fact, with a reason, when query quality or
scores are too weak to support an answer. \\
Verify mode & Require stronger support before treating a retrieved fact as
support for a proposed statement. \\
Local recall & Retrieve relevant Step or Fact chunks through the local embedding
path without requiring external embedding services. \\
Multilingual search & Recover relevant claims or Step evidence when the query
and stored material use different languages or mixed-language phrasing. \\
Time-aware search & Respect explicit time windows and temporal intent without
letting recency override stronger semantic or lexical evidence. \\
Local inference posture & Verify that \texttt{elephant-local-embed} can steady,
index, and retrieve at 64, 256, and 768 dimensions with the same Step/Fact
corpus. \\
Background reflect & Run a feature-scoped learning job and persist its summary
in \texttt{learning\_jobs.result\_json}. \\
\bottomrule
\end{longtable}
\endgroup

The evaluation question is whether Elephant Agent can continue the right work with the
right current understanding, and whether the user can correct that
understanding when it is wrong.

\subsection{Case Study Shapes}

The most informative case studies are multi-session and multi-surface:

\begin{itemize}[leftmargin=*]
  \item a repository refactor where Elephant Agent resumes with the correct Elephant context
  and recalls only relevant support;
  \item a design discussion where a correction to collaboration style changes
  the next reply;
  \item a dashboard fact correction where the old fact is retired and the new
  fact appears in the prompt;
  \item an answered Personal Model question that becomes an Identity, World,
  Pulse, or Journey fact with inspectable provenance;
  \item a recall query that correctly returns \texttt{no\_match} instead of
  inventing support;
  \item a background diary or Episode-close job that updates the Personal Model
  through tools and stores a learning result.
\end{itemize}

The reporting unit is the complete lifecycle: trigger, recalled support,
runtime action, fact update, provenance path, dashboard visibility, and later
wake behavior.
